\section*{Chapitre \ref{ch:g10:concentration-and-dilution} --- Concentration et dilution}

\begin{solution}{exo:g10:concentration-and-dilution:1}
$C_m = \qty{5.0}{g} / \qty{0.250}{L} = \qty{20}{g/L}$.
\end{solution}

\begin{solution}{exo:g10:concentration-and-dilution:2}
$c = \qty{0.050}{mol} / \qty{0.200}{L} = \qty{0.25}{mol/L}$.
\end{solution}

\begin{solution}{exo:g10:concentration-and-dilution:3}
$M(\ce{CuSO4}) = 63.5 + 32.1 + 4 \times 16.0 = \qty{159.6}{g/mol}$;
$n = 0.10 \times 0.1000 = \qty{0.0100}{mol}$; $m = 0.0100 \times 159.6 =
\qty{1.60}{g}$.
\end{solution}

\begin{solution}{exo:g10:concentration-and-dilution:4}
$F = 1.0 / 0.050 = 20$; $V_0 = \qty{100.0}{mL} / 20 = \qty{5.0}{mL}$.
\end{solution}

\begin{solution}{exo:g10:concentration-and-dilution:5}
$C_m = c \times M = 0.10 \times 58.5 = \qty{5.85}{g/L}$, soit environ
\qty{5.9}{g/L}.
\end{solution}

\begin{solution}{exo:g10:concentration-and-dilution:6}
$V_0 = c_1 V_1 / c_0 = 0.020 \times 250.0 / 0.50 = \qty{10.0}{mL}$
(\omterm{def:g10:concentration-and-dilution:dilution}{facteur de dilution} 25). Une pipette jaugée de \qty{10.0}{mL} et une
fiole jaugée de \qty{250.0}{mL}.
\end{solution}

\begin{solution}{exo:g10:concentration-and-dilution:7}
La fiole jaugée n'a qu'un trait, conçu pour un seul volume et précis à
une petite fraction de pour cent près~; les graduations d'un bécher ne
sont qu'une indication grossière. De même, une pipette jaugée délivre un
volume unique avec une grande précision, alors qu'une éprouvette graduée
est plus large et se lit moins précisément.
\end{solution}

\begin{solution}{exo:g10:concentration-and-dilution:8}
Entre \qty{0.02}{mol/L} et \qty{0.04}{mol/L}~: environ
\qty{0.03}{mol/L}. Pour plus de précision, ajouter des étalons entre ces
deux tubes, par exemple à 0.025, 0.030 et \qty{0.035}{mol/L}~: une
échelle plus fine là où il le faut.
\end{solution}

\begin{solution}{exo:g10:concentration-and-dilution:9}
$M(\ce{C6H12O6}) = \qty{180.0}{g/mol}$; $c = 50 / 180.0 =
\qty{0.28}{mol/L}$.
\end{solution}

\begin{solution}{exo:g10:concentration-and-dilution:10}
$M = \qty{342.0}{g/mol}$, donc $n = \qty{1.00}{mol}$ et $c =
\qty{1.00}{mol/L}$. Dilué dix fois~: $c = \qty{0.100}{mol/L}$, soit
$C_m = 0.100 \times 342.0 = \qty{34.2}{g/L}$. Une cuillerée de
\qty{0.020}{L} contient $34.2 \times 0.020 = \qty{0.68}{g}$ de
saccharose.
\end{solution}

\begin{solution}{exo:g10:concentration-and-dilution:11}
(a) La même masse de \omterm{def:g6:solutions-and-solubility:solute}{soluté} dans un volume deux fois plus grand~:
\qty{6.0}{g/L}. (b) Le volume est à peu près divisé par deux, la masse
de \omterm{def:g6:solutions-and-solubility:solute}{soluté} inchangée~: environ \qty{24}{g/L}.
\end{solution}

\begin{solution}{exo:g10:concentration-and-dilution:12}
\begin{enumerate}
  \item Trois \omterm{def:g10:concentration-and-dilution:dilution}{dilutions} par 10 font une \omterm{def:g10:concentration-and-dilution:dilution}{dilution} par $10^3$~:
        $\qty{0.10}{mol/L} / 1000 = \qty{1.0e-4}{mol/L}$.
  \item Une seule \omterm{def:g10:concentration-and-dilution:dilution}{dilution} par 1000 demanderait soit un volume minuscule
        de \omterm{def:g10:concentration-and-dilution:dilution}{solution mère} (\qty{0.1}{mL} dans \qty{100}{mL}), impossible à
        mesurer précisément avec une pipette, soit une très grande fiole
        (\qty{1}{mL} dans \qty{1}{L}). Trois \omterm{def:g10:concentration-and-dilution:dilution}{dilutions} au dixième
        utilisent des pipettes et des fioles ordinaires.
\end{enumerate}
\end{solution}

\begin{solution}{exo:g10:concentration-and-dilution:13}
\begin{enumerate}
  \item $M(\ce{CaCl2}) = 40.1 + 2 \times 35.5 = \qty{111.1}{g/mol}$;
        $n = 11.1 / 111.1 = \qty{0.100}{mol}$; $c = 0.100 / 0.5000 =
        \qty{0.200}{mol/L}$.
  \item Chaque \ce{CaCl2} donne un \ce{Ca^2+} et deux \ce{Cl-}~:
        $[\ce{Ca^2+}] = \qty{0.200}{mol/L}$ et $[\ce{Cl-}] =
        \qty{0.400}{mol/L}$.
\end{enumerate}
\end{solution}

\begin{solution}{exo:g10:concentration-and-dilution:14}
\begin{enumerate}
  \item $r = \qty{0.60}{cm}$, $h = \qty{0.20}{cm}$~:
        $V = \pi \times 0.60^2 \times 0.20 = \qty{0.23}{cm^3}$, soit
        \qty{0.23}{mL} d'eau en trop.
  \item Le même \omterm{def:g6:solutions-and-solubility:solute}{soluté} se trouve dans \qty{100.23}{mL} au lieu de
        \qty{100.00}{mL}~: la concentration est trop faible de
        $0.23 / 100.23 \approx
        \qty{0.23}{\percent}$.
\end{enumerate}
\end{solution}

\begin{solution}{exo:g10:concentration-and-dilution:15}
$n = 0.20 \times 0.100 + 0.10 \times 0.300 = 0.020 + 0.030 =
\qty{0.050}{mol}$ dans \qty{0.400}{L}~: $c = \qty{0.125}{mol/L}$. Ce
n'est pas la moyenne simple, \qty{0.15}{mol/L}~: il y a trois fois plus
de la solution diluée, si bien que le \omterm{def:g6:pure-substances-and-mixtures:mixture}{mélange} est plus proche de
\qty{0.10}{mol/L}. C'est la moyenne pondérée par les volumes.
\end{solution}

\begin{solution}{pb:g10:concentration-and-dilution:1}
\textbf{1.} \qty{0.9}{g} pour \qty{0.100}{L}~: $C_m = \qty{9.0}{g/L}$.

\textbf{2.} \qty{9.0}{g} pour \qty{1000}{mL}, soit \qty{9.0}{mg} par
millilitre.

\textbf{3.} $m = 9.0 \times 0.500 = \qty{4.5}{g}$.

\textbf{4.} Oui~: environ \qty{9}{g} par litre contre quelque
\qty{360}{g} par kilogramme d'eau à saturation, environ quarante fois
moins.

\textbf{5.} $9.0 \times 1.0 = \qty{9.0}{g}$ de sel.

\textbf{6.} $M(\ce{NaCl}) = 23.0 + 35.5 = \qty{58.5}{g/mol}$.

\textbf{7.} $c = 9.0 / 58.5 = \qty{0.154}{mol/L}$.

\textbf{8.} $n = 0.154 \times 0.500 = \qty{0.0769}{mol}$ (ou
$4.5 / 58.5$).

\textbf{9.} Chaque \ce{NaCl} donne un \ce{Na+} et un \ce{Cl-}~: tous deux
à \qty{0.154}{mol/L}.

\textbf{10.} $0.0769 \times \num{6.02e23} = \num{4.63e22}$ \omterm{def:g9:ions:ion}{ions} sodium.

\textbf{11.} $F = 200 / 9.0 = 22.2$.

\textbf{12.} Les \qty{4.5}{g} de la poche proviennent tous de la
solution concentrée~: $V_0 = \qty{4.5}{g} / \qty{200}{g/L} = \qty{0.0225}{L} =
\qty{22.5}{mL}$.

\textbf{13.} $V_0 = V_1 / F = 500 / 22.2 = \qty{22.5}{mL}$.

\textbf{14.} $c_0 = 200 / 58.5 = \qty{3.42}{mol/L}$~;
$c_0 V_0 = 3.42 \times 0.0225 = \qty{0.0769}{mol}$ et
$c_1 V_1 = 0.154 \times 0.500 = \qty{0.0769}{mol}$~: ils sont égaux.

\textbf{15.} Environ $500 - 22.5 = \qty{478}{mL}$ d'eau stérile.

\textbf{16.} Une pipette graduée de \qty{25}{mL} (ou une burette), lue
à \qty{0.1}{mL} près.

\textbf{17.} Dans une fiole jaugée de \qty{500.0}{mL}, complétée jusqu'au
trait de jauge.

\textbf{18.} $C_m = 200 \times 23.0 / 500 = \qty{9.20}{g/L}$, trop
concentrée de $0.20 / 9.0 \approx \qty{2.2}{\percent}$.

\textbf{19.} \qty{22.5}{mL} de la solution concentrée («~à 20~\%~») pour
une poche de \qty{500}{mL}.
\end{solution}
